\documentclass[12pt,a4paper]{article}
\usepackage[utf8]{inputenc}
\usepackage[T1]{fontenc}
\usepackage[spanish,es-noshorthands]{babel}
\usepackage{amsmath,amssymb,amsthm}
\usepackage{mathtools}
\usepackage{graphicx}
\usepackage{xcolor}
\usepackage{hyperref}
\usepackage{array,booktabs,multirow}
\usepackage[margin=2.5cm]{geometry}
\usepackage{setspace}
\onehalfspacing
\usepackage{verbatim}
\usepackage{listings}
\usepackage{caption}
\usepackage[numbers,sort&compress]{natbib}

\newtheorem{theorem}{Teorema}[section]
\newtheorem{lemma}[theorem]{Lema}
\newtheorem{proposition}[theorem]{Proposicion}
\newtheorem{corollary}[theorem]{Corolario}
\newtheorem{definition}[theorem]{Definicion}
\theoremstyle{remark}
\newtheorem{remark}{Observacion}[section]

\title{AViD Journal: verificación automatizada de novedad de teoremas\\
y los límites del enfoque}
\author{Ayrton Porto}
\date{Julio 2026}

\begin{document}
\maketitle

\begin{abstract}
Los sistemas de inteligencia artificial aplicados a la matemática verifican corrección pero no novedad: un teorema generado automáticamente puede compilar en Lean sin errores y ser, sin embargo, un resultado ya conocido. Este artículo presenta AViD Journal, un pipeline que recibe un artículo en LaTeX, formaliza sus enunciados en Lean 4 y emite un veredicto de novedad mediante un árbol de decisión sobre tres dimensiones: existencia previa en un corpus formal (Mathlib) y uno informal (TheoremSearch y Matlas, con filtro temporal y juez LLM), no-trivialidad por tácticas automáticas, y distancia estructural entre demostraciones medida como distancia de Jaccard sobre conjuntos de premisas.

La evaluación sobre artículos retirados de arXiv por duplicación declarada produjo un resultado más informativo que cualquier medida de desempeño: la identificación de tres obstáculos que limitan el enfoque con independencia de esta implementación. Primero, la compilación exitosa de un archivo Lean no garantiza fidelidad semántica. Segundo, el techo de recuperación lo impone la cobertura de los índices de teoremas y no la métrica de similitud. Tercero, arXiv elimina el código fuente de los artículos al retirarlos, lo que compromete la reproducibilidad de cualquier benchmark construido sobre ellos.
\end{abstract}

\section{Introducción}

El proyecto \textit{First Proof}~\cite{abouzaid2026}, en el que once matemáticos aportaron problemas de investigación no publicados, se diseñó para separar el razonamiento genuino de la mera búsqueda en la literatura. Su segundo lote, evaluado por jueces humanos~\cite{firstproof2b}, dejó documentado un modo de falla concreto de los modelos de lenguaje aplicados a la matemática: los sistemas rindieron mejor cuando el problema era estructuralmente similar a un resultado ya publicado, y en varios casos lo resolvieron traduciendo una demostración previa de la literatura ---reutilizando incluso su terminología y notación línea por línea--- sin citar la fuente, al punto de que un revisor humano lo habría marcado como plagio. Lo que falta, entonces, no es capacidad de generar demostraciones correctas, sino un método sistemático para evaluar si una demostración es genuinamente nueva. Los modelos pueden generar enunciados que compilan, que son correctos, y que sin embargo no son nuevos.

El problema es doble. Por un lado, los sistemas actuales de AI matemática (Axiom, Harmonic, DeepSeek-Prover) verifican corrección pero no novedad. Por otro, el criterio ingenuo que equipara novedad con ausencia en Mathlib~\cite{kasaura2025} falla en dos direcciones: marca como nuevos teoremas triviales que cualquier táctica automatizada de 2026 cierra sin necesidad de una idea matemática genuina, y omite resultados que no están formalizados pero sí publicados en la literatura informal.

Tao~\cite{tao2025}, en sus notas sobre \textit{Machine-Assisted Proof}, planteó la pregunta en términos de escala: si la generación automática de teoremas se acelera, hace falta filtrar novedad automáticamente. La verificación humana, que ya es el cuello de botella del peer review tradicional, se vuelve inviable cuando quien propone teoremas es un modelo.

La urgencia del problema quedó ilustrada mientras este trabajo se escribía. En
julio de 2026 se publicó un contraejemplo a la conjetura jacobiana, abierta
desde 1939, encontrado con asistencia de un modelo de lenguaje y verificable de
forma elemental~\cite{alpoge2026}. Como la conjetura contaba con numerosos
resultados equivalentes conocidos, su refutación arrastró automáticamente a
varios de ellos, y en los días siguientes comenzaron a aparecer en arXiv
contraejemplos derivados a otras conjeturas. Un episodio así produce de forma
concentrada el material que motiva este trabajo: una cantidad de resultados que
descienden de un mismo hallazgo, obtenidos en paralelo por autores distintos,
donde la pregunta de si un enunciado ya fue establecido por otro es genuina y
costosa de responder manualmente. Es además el régimen en que la limitación de
cobertura que documentamos en la sección~\ref{ssec:cobertura} no se aplica: los
trabajos potencialmente duplicados son preprints recientes, indexados desde su
aparición.

Este artículo presenta AViD Journal (Automated Verification in Demonstrations), un sistema que aborda esa pregunta de manera operacional: recibe un artículo en LaTeX, formaliza sus enunciados en Lean 4, y emite un veredicto de novedad mediante un árbol de decisión en tres dimensiones. Construir el camino completo permitió una segunda contribución, que resultó ser la más informativa: identificar dónde falla el enfoque. El eslabón débil no es la búsqueda ni el filtro de trivialidad, que funcionan; es que la verificación por compilación no garantiza fidelidad semántica de la formalización, y que la cobertura de los corpus de recuperación disponibles no alcanza a los resultados que la duplicación matemática típicamente re-descubre. Este artículo documenta ambos muros con la misma atención que el pipeline.

\subsection{Contribuciones}

\begin{enumerate}
\item \textbf{Un pipeline completo de verificación de novedad.} El sistema recorre el camino entero, desde el fuente LaTeX hasta el veredicto: parseo de entornos matemáticos, formalización en Lean 4 con abstracción multi-modelo, y un árbol de decisión de ocho veredictos sobre tres dimensiones. D1 verifica existencia previa en un corpus formal (Mathlib, vía Leandex) e informal (TheoremSearch y Matlas, con filtro temporal y un juez LLM); D2 filtra trivialidad mediante tácticas automáticas de Lean; D3 mide distancia estructural entre demostraciones como distancia de Jaccard sobre conjuntos de premisas.

\item \textbf{Un benchmark de papers retirados por duplicación.} Se construyó un dataset de papers retirados de arXiv cuyo comentario de retiro declara duplicación de resultados previos, con controles emparejados por categoría y año, junto con el análisis de la recuperabilidad de sus duplicadores.

\item \textbf{Un mapa de los modos de falla del enfoque.} La contribución central de la evaluación no es una métrica de desempeño sino la caracterización de tres muros: (a) la compilación exitosa de un archivo Lean no garantiza que la formalización sea fiel al enunciado original, y documentamos cinco modos de falla sucesivos que ilustran por qué ninguna guardia sintáctica cierra el problema; (b) los duplicadores que la duplicación matemática re-descubre son mayoritariamente anteriores a la era de los preprints electrónicos, y por lo tanto invisibles para los corpus de recuperación disponibles; (c) arXiv elimina la fuente LaTeX de los papers al retirarlos, lo que impide procesarlos desde su código fuente y limita la reproducibilidad de cualquier benchmark construido sobre ellos.
\end{enumerate}

\subsection{Estructura del artículo}

La sección~\ref{sec:related} sitúa este trabajo en la intersección de cuatro frentes de investigación. La sección~\ref{sec:pipeline} describe el pipeline: parser, formalización, corpus de referencia, las tres dimensiones y el árbol de veredictos. La sección~\ref{sec:instrumentos} valida cada instrumento por separado. La sección~\ref{sec:experimentos} reporta la evaluación sobre papers retirados y caracteriza los modos de falla encontrados. La sección~\ref{sec:limitaciones} enumera las limitaciones del framework, la implementación y el diseño experimental. La sección~\ref{sec:conclusion} cierra con lo construido y lo pendiente.
\section{Trabajo relacionado}
\label{sec:related}

Determinar si un teorema es nuevo toca tres frentes que rara vez se cruzan: la infraestructura de búsqueda de enunciados, los filtros de novedad en generación de conjeturas, y la identidad de demostraciones.

\subsection{Buscadores e infraestructura de teoremas}

\textbf{TheoremSearch}~\cite{theoremsearch2026} indexa 9.2 millones de enunciados extraídos de arXiv y siete fuentes adicionales, representa cada teorema con una descripción corta en lenguaje natural para embeber, y expone una API pública sin autenticación. Su motivación documentada —los retiros por duplicación en arXiv— es la misma que motiva nuestro experimento con papers retirados. TheoremSearch encuentra enunciados similares; AViD agrega la capa de veredicto.

\textbf{Matlas}~\cite{matlas2026} extrae 8.07 millones de enunciados de 435.000 artículos revisados por pares que abarcan de 1826 a 2025, provenientes de 180 revistas seleccionadas por criterio de citación ICM, más 1.900 libros de texto. Las dos fuentes son complementarias en cobertura temporal: TheoremSearch cubre arXiv desde 1991, Matlas cubre revistas peer-reviewed desde 1826. AViD las consume como las dos ramas de su corpus informal, y la sección~\ref{sec:experimentos} reporta hasta dónde alcanza esa cobertura combinada.

El espacio de buscadores sobre Mathlib está saturado. \textbf{LeanExplore}~\cite{leanexplore2025}, backend de Leandex, ofrece búsqueda sobre declaraciones de Lean 4; \textbf{Lean Finder}~\cite{leanfinder2025} incorpora intención del usuario a la búsqueda semántica sobre Mathlib. Existen además herramientas de búsqueda por patrón y por estado de prueba [VERIFICAR: Loogle, LeanSearch, LeanStateSearch]. AViD no compite en ese espacio: consume una de ellas como infraestructura (Leandex es el backend de D1 C\_F) y agrega la capa de decisión que ninguna ofrece. Un survey reciente sobre IA matemática~\cite{survey2026} cita a los buscadores de enunciados como infraestructura para determinar si un resultado ya es conocido, lo que sitúa el problema que AViD aborda como reconocido por el campo.

% [VERIFICAR] TheoremGraph + LeanGraph (arXiv:2606.25363): grafo unificado
% formal-informal, 11.7M entornos tipo-teorema, LeanGraph 388.105 nodos.
% Si verifica, va acá como competidor más cercano en infraestructura.
% [VERIFICAR] COMPOSE (arXiv:2605.30333): predice teoremas futuros desde
% estructura de citas. Si verifica, va acá como dirección complementaria.

\subsection{Filtros de novedad en generación de conjeturas}

\textbf{Kasaura et al.}~\cite{kasaura2025} [VERIFICAR arXiv:2509.14274] definen novedad como ausencia en Mathlib, en la librería generada durante la sesión, y en una lista manual del conjeturador. Es el baseline que AViD corrige: sin filtro de trivialidad (D2), un teorema como ``la suma de cuatro pares es par'' se marcaría como nuevo; sin distancia estructural (D3), dos pruebas distintas del mismo enunciado serían indistinguibles.

\textbf{LeanConjecturer}~\cite{leanconjecturer2025} [VERIFICAR arXiv:2506.22005] filtra novedad con \texttt{exact?} contra Mathlib y no-trivialidad con \texttt{aesop}: generó 12.289 conjeturas desde 40 archivos semilla, de las cuales 3.776 resultaron sintácticamente válidas y no triviales. Son exactamente los mecanismos que AViD implementa como D1 (\texttt{exact?} como fallback de C\_F) y D2 (\texttt{aesop} en \texttt{T\_AUTO}). La contribución de AViD no es inventar estos filtros sino componerlos en un árbol de decisión que emite un veredicto, agregar D3 y una rama informal con filtro temporal, y evaluarlos contra ground truth externo. Trabajos sobre pruning de conjeturas redundantes~\cite{mining2024} [VERIFICAR arXiv:2412.16177] y generación sintética por forward-reasoning~\cite{synthetic2024} [VERIFICAR OpenReview EeDSMy5Ruj] ilustran la misma necesidad desde el lado de la generación.

\subsection{Identidad y similitud de demostraciones}

La distancia de Jaccard sobre conjuntos de premisas que AViD usa en D3 tiene genealogía en la literatura de premise selection. La línea MaSh~\cite{kuhlwein2013} introdujo aprendizaje automático para la selección de premisas en Sledgehammer, y \textbf{Kaliszyk y Urban}~\cite{kaliszyk2014} usaron k-NN con features ponderadas por IDF sobre el corpus Flyspeck; \textbf{Magnushammer}~\cite{magnushammer2024} aplica la misma tarea con transformers. AViD reposiciona la herramienta: en lugar de usarla como entrada para construir pruebas, la usa como huella para comparar pruebas ya construidas. \textbf{Piotrowski et al.}~\cite{piotrowski2023} [VERIFICAR arXiv:2304.00994] introducen un \textit{math filter} que preserva solo lemas de naturaleza claramente matemática usando nombres de Mathlib como whitelist; los filtros previos al Jaccard de D3 siguen esa receta.

El análisis de red de Mathlib de Li et al.~\cite{network2026} sobre 308.129 declaraciones y 8.4 millones de aristas encuentra que la centralidad de red captura infraestructura del lenguaje más que relevancia matemática. Ese hallazgo respalda el Filtro 1 de D3: al eliminar premisas de los namespaces \texttt{Init.} y \texttt{Lean.} antes de calcular Jaccard, AViD no aplica una heurística ad-hoc sino un filtro consistente con evidencia publicada.

\textbf{Yoo}~\cite{yoo2025} [VERIFICAR arXiv:2504.00063] representa teoremas como \textit{proof vectors} sobre sistemas de axiomas y los compara con similitud coseno, euclidiana o Jaccard. Es el trabajo más cercano en herramienta conceptual pero con propósito opuesto: el Atlas organiza teoremas por similitud estructural; AViD usa esa similitud para decidir si una prueba es redundante. \textbf{Huch}~\cite{huch2022} [VERIFICAR arXiv:2209.13305] documenta la distribución scale-free del grado de entrada en el AFP, relevante para el trabajo futuro con premisas ponderadas por IDF.

\subsection{Corrección frente a novedad}

Los pipelines paper-to-Lean existentes verifican que una demostración sea correcta, no que sea nueva; la verificación de novedad queda delegada a revisores humanos. AViD es explícitamente esa capa delegada. Un sistema completo de revisión automatizada requeriría ambas dimensiones, y la sección~\ref{sec:experimentos} muestra que la primera no es condición suficiente para la segunda: un archivo Lean puede compilar sin errores y no ser una formalización fiel del enunciado original.

% [VERIFICAR] MerLean: pipeline paper-to-Lean que delega novedad a humanos.
% [VERIFICAR] Pseudo-Formalization / ArxivMathGradingBench: verifica corrección
% de pruebas de arXiv vía formalización. Vecino exacto en el espacio de diseño.
% Si verifican, esta subsección los nombra en lugar de la formulación genérica.

\subsection{Posicionamiento}

\begin{table}[htbp]
\centering
\small
\setlength{\tabcolsep}{4pt}
\begin{tabular}{lccccc}
\toprule
Sistema & Existencia & Trivialidad & Dist.\ prueba & Corpus informal & Veredicto \\
\midrule
Kasaura et al.~\cite{kasaura2025}      & Mathlib            & --- & --- & --- & sí/no \\
LeanConjecturer~\cite{leanconjecturer2025} & \texttt{exact?} & \texttt{aesop} & --- & --- & --- \\
TheoremSearch~\cite{theoremsearch2026} & 9.2M enunciados    & --- & --- & \checkmark & --- \\
Matlas~\cite{matlas2026}               & 8.07M enunciados   & --- & --- & \checkmark & --- \\
Atlas~\cite{yoo2025}                   & ---                & --- & Jaccard & --- & --- \\
AViD (este trabajo)                    & Mathlib + informal & 6 tácticas & Jaccard$^{*}$ & TS + Matlas & 8 veredictos \\
\bottomrule
\end{tabular}
\caption{Posicionamiento de AViD respecto de trabajos previos. Los trabajos previos construyen piezas de infraestructura; AViD las compone en un pipeline que va del fuente LaTeX al veredicto. $^{*}$D3 está implementada y validada de forma aislada (sección~\ref{sec:instrumentos}), pero no se ejercita dentro del árbol de decisión en los experimentos reportados.}
\label{tab:posicionamiento}
\end{table}

La diferencia no está en ninguna celda individual sino en la composición: los trabajos previos construyen piezas de infraestructura —búsqueda, formalización, similitud—; AViD las ensambla en un pipeline que recibe un paper en LaTeX y devuelve un veredicto.

La diferencia no está en ninguna celda individual sino en la composición: los trabajos previos construyen piezas de infraestructura —búsqueda, formalización, similitud—; AViD las ensambla en un pipeline que recibe un paper en LaTeX y devuelve un veredicto.

\section{El pipeline AViD}
\label{sec:pipeline}

El sistema recibe un archivo \texttt{.tex}, extrae sus bloques matemáticos, los formaliza en Lean 4, y aplica un árbol de decisión en tres dimensiones para emitir un veredicto de novedad. Esta sección describe cada etapa con el detalle necesario para su reproducción conceptual.

\subsection{Ingesta y parseo del LaTeX}

El parser extrae entornos matemáticos del fuente LaTeX: \texttt{theorem}, \texttt{lemma}, \texttt{proposition}, \texttt{corollary}, \texttt{definition}, y sus variantes en español (\texttt{teorema}, \texttt{lema}, \texttt{proposición}, \texttt{definición}, \texttt{corolario}). También reconoce variantes abreviadas (\texttt{thm}, \texttt{lem}, \texttt{prop}, \texttt{cor}, \texttt{defn}) y entornos definidos por el usuario mediante \verb|\newtheorem| y \verb|\theoremstyle|.

Las dependencias entre bloques se extraen de las referencias cruzadas \verb|\ref{}| y \verb|\cite{}|. Con ellas el parser construye un grafo dirigido acíclico y aplica un ordenamiento topológico (algoritmo de Kahn) para determinar la secuencia de formalización: los bloques sin dependencias se procesan primero; los que dependen de otros, después.

No todos los entornos detectados pasan a la etapa siguiente. El orquestador solo formaliza tipos considerados \textit{formalizables}: \texttt{theorem}, \texttt{lemma}, \texttt{proposition}, \texttt{corollary} y sus variantes. Entornos como \texttt{remark}, \texttt{example}, \texttt{proof} y \texttt{definition} (cuando no es requerida como dependencia) se extraen pero no se formalizan.

\textbf{Limitación conocida.} El parser opera por expresión regular compilada con una lista fija de nombres de entorno. Los papers que usan AMS-\TeX{} (\verb|\documentstyle{amsppt}|) o definen entornos con nombres idiosincráticos (\verb|\begin{inizio}|, \verb|\begin{numero}|) no son parseables por esta versión. De los 33 candidatos del dataset de retirados, 7 resultaron no viables por este motivo.

\subsection{Formalización a Lean 4}

\subsubsection{Abstracción multi-modelo}

La formalización de cada bloque (enunciado y demostración) se delega a un modelo de lenguaje a través de una interfaz abstracta \texttt{ModelProvider}. La implementación actual soporta ocho backends: Claude Code (modo agéntico, vía CLI), Anthropic API, OpenAI, DeepSeek, OpenRouter, Mistral, Gemini, y OpenCode Go. Los experimentos reportados en la sección~\ref{sec:experimentos} usan Qwen 3.7-max a través de OpenCode Go, seleccionado por comparación empírica entre backends.

La arquitectura distingue dos familias de proveedores. Los \textit{agénticos} manejan su propio ciclo de verificación: el modelo recibe el prompt, genera código Lean, intenta compilar, y corrige iterativamente. Los \textit{API} usan un ciclo de verificación externo: el sistema envía el prompt, recibe la respuesta, extrae el código Lean, compila con \texttt{lake env lean}, y reenvía los errores como feedback para la siguiente iteración.

Cada bloque se clasifica en uno de cuatro modos según su complejidad estimada: \texttt{SIMPLE} (5 rondas máximas de corrección), \texttt{MEDIUM} (15 rondas), \texttt{HARD} (30 rondas), y \texttt{EXTERNAL} (sin modelo; se emite un axioma con referencia a la fuente).

\subsubsection{Proyecto Lean compartido}

Todos los papers comparten un único proyecto Lean 4 (v4.29.0) con una compilación de Mathlib (8.247 archivos \texttt{.olean}). Cada paper se aloja como submódulo en \texttt{Papers/<Slug>/}. Los bloques formalizados se acumulan incrementalmente en un archivo \texttt{Paper.lean} y se registran en un índice que el orquestador consulta antes de buscar en Mathlib.

El criterio de éxito de formalización es exigente: el archivo debe compilar sin errores, sin \texttt{sorry} en ninguna declaración, y debe contener al menos una declaración sustantiva (\texttt{theorem}, \texttt{lemma}, \texttt{def}). Los archivos vacíos o con únicamente imports se rechazan. La sección~\ref{sec:experimentos} documenta por qué este criterio, aun siendo exigente, no garantiza fidelidad semántica.

\subsubsection{Modo statement-only}

Para experimentos donde solo se requiere el enunciado, el sistema opera en modo \textit{statement-only}: el modelo genera el enunciado en Lean con \texttt{:= by sorry} y el criterio de éxito se relaja para aceptar la presencia de \texttt{sorry} siempre que no haya errores de compilación. Este modo no es una capacidad nativa del código base sino una configuración de los prompts y del criterio de aceptación. Todos los experimentos sobre papers reales reportados en este trabajo operan en este modo, lo que implica que D3 —que requiere premisas de una demostración— no se ejercita en ellos.

\subsection{Corpus de referencia}

El sistema compara cada enunciado contra tres corpus.

\textbf{C\_F (corpus formal).} Mathlib v4.29.0, accedido a través de la API de Leandex, cuyo backend es LeanExplore~\cite{leanexplore2025}. Leandex indexa los enunciados de Mathlib y devuelve matches con su estatus de prueba. La versión actual de la API (v2) no proporciona scores de similitud; el sistema usa el ordenamiento provisto por la API y descarta resultados cuyo estatus es \texttt{statement\_only}. La ausencia de umbral es una limitación con consecuencias sobre la interpretación de los resultados de D1 C\_F (sección~\ref{sec:instrumentos}).

\textbf{C\_I (corpus informal).} Dos fuentes complementarias en cobertura temporal: TheoremSearch~\cite{theoremsearch2026}, con 9.2 millones de enunciados extraídos de arXiv y siete fuentes adicionales, y Matlas~\cite{matlas2026}, con 8.07 millones de enunciados de revistas peer-reviewed desde 1826. Ambas exponen API pública sin autenticación. Los candidatos recuperados se filtran con embeddings MiniLM (umbral de similitud de coseno $\geq 0.40$) y luego pasan por los filtros descritos en la sección~\ref{ssec:d1} antes de llegar al juez LLM.

\textbf{Corpus propio.} A medida que el orquestador formaliza bloques de un paper, los acumula en un índice local. Los bloques ya procesados dentro del mismo paper se consultan antes de buscar en Mathlib, evitando que el sistema marque como ya existente un teorema que el propio paper acaba de introducir.

\subsection{Las tres dimensiones}

\subsubsection{D1: no-existencia previa}
\label{ssec:d1}

D1 verifica si el enunciado ya aparece en el corpus formal (C\_F) o informal (C\_I). La evaluación sigue un orden fijo con cortocircuito: si C\_F encuentra match, C\_I no se ejecuta.

\textbf{Rama C\_F.} Se consulta la API de Leandex con el enunciado en lenguaje natural. Si Leandex devuelve al menos un resultado cuyo estatus no sea \texttt{statement\_only}, se considera que el teorema existe en Mathlib. No se aplica umbral de similitud, porque la API no proporciona scores.

\textbf{Fallback \texttt{exact?}.} Si Leandex no encuentra match, el sistema ejecuta \texttt{example : $\tau$ := by exact?} con un presupuesto de 15 segundos sobre el entorno Lean con Mathlib completo. Si la táctica encuentra un teorema que cierra el goal, se trata como match en C\_F. Esta táctica se clasifica como D1 y no como D2 porque su función es encontrar un teorema ya demostrado, no cerrar el goal por automatización.

\textbf{Rama C\_I.} Solo se activa si C\_F no encontró match. La etapa A consulta TheoremSearch y Matlas con el enunciado, previamente limpiado de marcado LaTeX. Los candidatos supervivientes al umbral de embeddings pasan por dos filtros estructurales antes de la etapa B, en la que un juez LLM (temperature $=0$) compara cada candidato con el enunciado original y emite uno de cuatro veredictos: \texttt{equivalent}, \texttt{generalization}, \texttt{specialization}, \texttt{different}.

\textbf{Filtro temporal.} Un duplicador válido debe ser anterior al trabajo que duplica. El sistema descarta todo candidato cuya fecha de publicación sea posterior a la del paper bajo evaluación. Para candidatos de arXiv, la fecha se deriva del identificador, que codifica año y mes en tres formatos históricos distintos (\texttt{YYMM.NNNNN}, \texttt{cat/YYMMNNN}, y la variante sin prefijo de categoría). Para candidatos de Matlas, que se identifican por DOI y no por arXiv ID, se usa el año de publicación provisto por la fuente. Los candidatos sin fecha recuperable no se descartan: pasan a la etapa B y se contabilizan por separado, de modo que su presencia quede registrada en lugar de producir un descarte silencioso. Si el paper bajo evaluación no tiene fecha parseable, el filtro se desactiva.

Sin este filtro, el sistema puede emitir un veredicto de no-novedad contra un trabajo \textit{posterior} al evaluado, lo cual es lógicamente imposible como duplicación. La sección~\ref{sec:experimentos} documenta un caso concreto y su corrección.

\textbf{Exclusión del propio paper.} Los índices de C\_I contienen los papers que el sistema evalúa. Sin exclusión explícita, un paper se recupera a sí mismo como candidato y el juez lo clasifica como \texttt{equivalent}, produciendo un falso positivo de duplicación. El sistema pasa el identificador del paper evaluado como parámetro de exclusión a las fuentes que lo soportan de forma nativa, y aplica un filtro posterior a la búsqueda sobre las que no.

\subsubsection{D2: no-trivialidad}

D2 verifica si el enunciado puede demostrarse exclusivamente con tácticas automáticas estándar de Lean. Si alguna táctica cierra \texttt{example : $\tau$ := by T}, el teorema se considera trivial.

\textbf{Tácticas, orden y presupuestos.} El conjunto \texttt{T\_AUTO} se ejecuta en este orden: \texttt{decide}, \texttt{norm\_num}, \texttt{simp}, \texttt{omega}, \texttt{tauto}, \texttt{aesop}. Las primeras cinco disponen de 10 segundos cada una; \texttt{aesop} dispone de 30. El orden prioriza tácticas baratas y específicas; \texttt{aesop}, que realiza una búsqueda más exhaustiva, va al final. La ejecución se detiene en la primera táctica que cierra el goal.

% [VERIFICAR/AJUSTAR] El texto original afirmaba un overhead de 45s por
% invocación de lake env lean, incompatible con los tiempos de 13s reportados
% en la Tabla de la seccion 4.2. Reformulado abajo asumiendo precalentamiento
% de oleans. AJUSTAR a lo que el instrumento mide realmente.
Cada invocación de táctica requiere cargar el entorno Lean con Mathlib completo. Sin precalentamiento, ese arranque domina el tiempo de ejecución; los tiempos reportados en la sección~\ref{sec:instrumentos} corresponden a invocaciones sobre un entorno ya cargado.

\textbf{Blacklist de \texttt{norm\_num}.} La táctica \texttt{norm\_num} en Mathlib v4.29.0 incluye un atajo que cierra \texttt{Irrational (Real.sqrt 2)}. Para evitar ese falso positivo, \texttt{norm\_num} se excluye cuando el enunciado contiene la palabra \texttt{Irrational}.

\textbf{\texttt{exact?} no pertenece a D2.} La táctica fue reubicada en D1 como fallback de C\_F por una razón semántica: busca un teorema existente en el entorno, lo cual es verificación de existencia previa y no de trivialidad.

\textbf{Sensibilidad al enunciado formalizado.} El resultado de D2 no es una propiedad fija del teorema sino del par (enunciado formalizado, táctica, presupuesto). Un mismo teorema informal puede producir formalizaciones que \texttt{aesop} cierra en segundos o que exceden el presupuesto por un orden de magnitud, según qué definiciones haya elegido el formalizador. La sección~\ref{sec:instrumentos} documenta un caso concreto.

\subsubsection{D3: distancia estructural de premisas}

D3 mide cuán diferente es una demostración de otra existente en Mathlib. La métrica es la distancia de Jaccard sobre los conjuntos de premisas (lemas, teoremas, definiciones) usados en cada prueba.

\textbf{Extracción de premisas.} Las premisas se extraen con una herramienta standalone en Lean que recorre el \texttt{InfoTree} y recolecta, para cada nodo \texttt{TermInfo}, la constante referenciada, su ubicación de definición y el módulo que la contiene. Se excluyen las referencias que ocurren en el sitio mismo de definición de la constante, para evitar que un teorema se registre a sí mismo como premisa. El extractor funciona en Windows nativo, sin requerir WSL ni LeanDojo; un wrapper en Python gestiona la ejecución y mantiene un caché por hash del archivo fuente.

\textbf{Identidad canónica y deduplicación.} Dos premisas con la misma ubicación de definición representan el mismo objeto lógico, aunque aparezcan en posiciones distintas de la prueba. Antes de cualquier filtro, las premisas se deduplican por esta identidad canónica.

\textbf{Pipeline de filtrado.} Antes de calcular Jaccard, las premisas pasan por dos filtros:

\begin{enumerate}
\item \textbf{Namespace blacklist.} Se eliminan las premisas cuyo módulo comienza con \texttt{Init.} o \texttt{Lean.}. Estos namespaces contienen constructores de tipo del núcleo, instancias de typeclasses e internals de tácticas que el elaborador resuelve automáticamente; no reflejan contenido matemático y dominarían artificialmente la intersección. El análisis de red de Mathlib de Li et al.~\cite{network2026} respalda este criterio: sobre un grafo de 308.129 declaraciones y 8,4 millones de aristas, encuentran que la centralidad captura infraestructura del lenguaje más que relevancia matemática.

\item \textbf{Premisas del enunciado.} Se eliminan las premisas cuya posición en el archivo fuente cae dentro del rango de líneas del enunciado. Las constantes que aparecen en el enunciado (hipótesis, tipos, definiciones locales) son parte de la firma y no de la prueba; compararlas infla artificialmente la similitud.
\end{enumerate}

\textbf{Cálculo.} Sean $P(A)$ y $P(B)$ los conjuntos de premisas post-filtros y deduplicados. La distancia de Jaccard es
\[
d_J(A,B) \;=\; 1 - \frac{|P(A) \cap P(B)|}{|P(A) \cup P(B)|}.
\]
El umbral de decisión es $\theta = 0{,}5$: si la distancia lo supera, las pruebas se consideran estructuralmente distintas. Si alguno de los conjuntos queda vacío tras los filtros, no se emite distancia y el veredicto es \texttt{INCONCLUSIVE}.

\textbf{Estado de calibración.} El umbral $\theta = 0{,}5$ es el valor inicial de diseño y no ha sido calibrado contra un conjunto amplio de pares. La sección~\ref{sec:instrumentos} reporta su comportamiento sobre cinco pares con juicio humano, de los cuales solo dos resultan informativos para calibración.

\subsection{Árbol de veredictos}

El orquestador recorre las tres dimensiones en orden de costo creciente (figura~\ref{fig:arbol}).

\begin{figure}[htbp]
\centering
\small
\begin{verbatim}
D2 (trivialidad): alguna tactica de T_AUTO cierra tau?
  |
  +-- SI  -> NO_NOVEDOSO_trivial                      (FIN)
  |
  +-- NO  -> D1 C_F (Leandex)
             |
             +-- match -> D3 (si hay premisas)
             |            +-- distantes      -> NOVEDAD_DEMOSTRACION
             |            +-- cercanas       -> NO_NOVEDOSO_redundante
             |            +-- conjuntos vacios -> INCONCLUSIVE
             |            +-- no disponible  -> MATCH_PENDIENTE_D3
             |
             +-- sin match -> exact? (fallback)
                              |
                              +-- match -> mismo camino que C_F
                              |
                              +-- sin match -> D1 C_I
                                               (TheoremSearch + Matlas,
                                                filtro temporal, juez LLM)
                                   +-- equivalent        -> CONOCIDO_LITERATURA
                                   +-- generalization /
                                       specialization    -> ZONA_GRIS
                                   +-- different o
                                       sin candidatos    -> NOVEDAD_ENUNCIADO
\end{verbatim}
\caption{Árbol de decisión del orquestador. Las dimensiones se evalúan en orden de costo creciente: D2 es local, D1 C\_F es una consulta HTTP, D1 C\_I involucra varias APIs y un juez LLM, y D3 requiere extracción de premisas.}
\label{fig:arbol}
\end{figure}

\textbf{Lógica de cortocircuito.} D2 se evalúa primero porque es local y barata: si el teorema es trivial, el árbol termina sin consultar APIs externas. D1 C\_F es la segunda por ser una consulta HTTP rápida. D1 C\_I es la tercera porque involucra varias APIs y una llamada al juez. D3 es la más cara y solo se ejecuta cuando hay match en C\_F que requiere dirimir si la prueba es novedosa o redundante.

\textbf{Los ocho veredictos.}

\begin{enumerate}
\item \texttt{NOVEDAD\_ENUNCIADO}: sin match en C\_F ni C\_I.
\item \texttt{NOVEDAD\_DEMOSTRACION}: match en C\_F, D3 indica pruebas distantes.
\item \texttt{CONOCIDO\_LITERATURA}: match en C\_I, sin match en C\_F.
\item \texttt{NO\_NOVEDOSO\_redundante}: match en C\_F, D3 indica misma prueba.
\item \texttt{NO\_NOVEDOSO\_trivial}: D2 cierra con táctica estándar.
\item \texttt{ZONA\_GRIS}: el juez clasifica el match como generalización o especialización; requiere revisión humana.
\item \texttt{MATCH\_PENDIENTE\_D3}: match en C\_F, D3 no disponible.
\item \texttt{INCONCLUSIVE}: D3 ejecutado, conjuntos de premisas vacíos tras los filtros.
\end{enumerate}

Tres de los ocho veredictos (\texttt{NOVEDAD\_DEMOSTRACION}, \texttt{NO\_NOVEDOSO\_redundante}, \texttt{INCONCLUSIVE}) dependen de D3 y, por lo tanto, no se emiten en los experimentos reportados, que operan en modo statement-only.

\subsection{Implementación y estado actual}

El sistema está implementado en Python 3.11+ y corre sobre Windows 10 nativo; toda la infraestructura, incluida la extracción de premisas, funciona sin requerir WSL ni Docker. La base de código se organiza en dos módulos: uno contiene el parser, el orquestador de novedad y las tres dimensiones; el otro, un orquestador de formalización con abstracción multi-modelo. El árbol de veredictos es idéntico independientemente de qué orquestador haya formalizado el paper.

El proyecto cuenta con 153 pruebas superadas y 1 saltada, un conjunto de evaluación de 24 teoremas, y un dataset de 26 papers retirados por duplicación con controles emparejados por categoría y año.


\section{Validación de instrumentos}
\label{sec:instrumentos}

Antes de evaluar el sistema contra papers reales (sección~\ref{sec:experimentos}), esta sección valida cada instrumento por separado contra ground truth conocido: D3 contra pares de calibración con juicio humano, D2 contra el conjunto de evaluación, y D1 C\_F contra la cobertura del mismo conjunto.

\subsection{D3: escalera de calibración}
\label{ssec:d3cal}

La distancia de Jaccard se validó sobre cinco pares de teoremas para los cuales un juez humano (el primer autor) determinó la relación esperada entre sus demostraciones y firmó cada etiqueta antes de ejecutar la medición. Los pares cubren el espectro: misma prueba (control \textit{self}), pruebas genuinamente distintas, pruebas aparentemente distintas que resultaron idénticas tras la formalización, y teoremas no relacionados (control cruzado).

\begin{table}[htbp]
\centering
\small
\setlength{\tabcolsep}{5pt}
\begin{tabular}{llrrrr}
\toprule
Par & Etiqueta del juez & Jaccard & Distancia & $|\cap|$ & $|\cup|$ \\
\midrule
T08a vs.\ T08a & control: misma prueba      & 1{,}0000 & 0{,}0000 & 9 & 9 \\
T07a vs.\ T07b & \textit{same disguised}    & 0{,}5000 & 0{,}5000 & 1 & 2 \\
T08a vs.\ T08b & \textit{genuinely different} & 0{,}2778 & 0{,}7222 & 5 & 18 \\
T09a vs.\ T09b & \textit{genuinely different} & 0{,}0000 & 1{,}0000 & 0 & 6 \\
T07 vs.\ T08   & control: no relacionados   & 0{,}0000 & 1{,}0000 & 0 & 10 \\
\bottomrule
\end{tabular}
\caption{Calibración de D3 sobre cinco pares con etiqueta humana previa. Los conjuntos de premisas son post-filtros y deduplicados.}
\label{tab:d3cal}
\end{table}

\textbf{T08 (distancia 0{,}7222).} Dos pruebas genuinamente distintas de la irracionalidad de $\sqrt{2}$ producen distancia alta. Las cinco premisas compartidas son lemas fundacionales de aritmética; las restantes, exclusivas de cada lado, capturan las estrategias divergentes (divisibilidad por primos frente a valuación 2-ádica). El umbral $\theta = 0{,}5$ clasifica correctamente el par como pruebas distantes.

\textbf{T09 (distancia 1{,}0).} Las dos pruebas de la suma de Gauss —inducción con 

\noindent \texttt{sum\_range\_succ} frente a fórmula cerrada con \texttt{Finset.sum\_range\_id}— no comparten ninguna premisa tras los filtros. Cada lado aporta seis premisas exclusivas de su estrategia, y la distancia máxima refleja correctamente la diferencia etiquetada por el juez.

\textbf{T07 (distancia 0{,}50): punto degenerado.} Las dos pruebas de la infinitud de los primos (Euclides y Euler) quedan exactamente sobre el umbral. Tras la formalización, ambas colapsaron a la misma invocación de \texttt{Nat.exists\_infinite\_primes}. Con solo dos premisas totales tras los filtros y una compartida, el Jaccard es frágil: cambios menores en el comportamiento de los filtros invertirían el veredicto. El juez las había etiquetado como \textit{same disguised}: enunciados distintos en el fuente LaTeX, idénticos en el término de prueba elaborado. Este par no es informativo para calibración.

\textbf{Control negativo.} La distancia entre un teorema de teoría de números y uno de análisis es 1{,}0 con intersección vacía, lo que confirma que D3 no produce similitud espuria entre dominios distintos.

\textbf{Qué calibra esta escalera.} Los cinco pares son consistentes con el
juicio humano bajo $\theta = 0{,}5$. Sin embargo, tres de ellos (T09 y los dos
controles) caen en los extremos del rango y no restringen la elección del
umbral: cualquier valor en $(0,1)$ los clasifica correctamente. Las únicas
cotas efectivas provienen de T08, que exige $\theta < 0{,}7222$, y de T07, que
exige $\theta \geq 0{,}50$. Como T07 es degenerado ---unión de dos premisas,
intersección de una, de modo que un cambio menor en el comportamiento de los
filtros invertiría el veredicto---, la cota inferior no es confiable. El
intervalo admisible sostenido por evidencia se reduce entonces a
$\theta < 0{,}7222$, y el valor $\theta = 0{,}5$ es una elección de diseño
dentro de ese rango, no un óptimo calibrado. Estrechar ese intervalo requiere
pares cuya distancia caiga en el interior del rango y cuyos conjuntos de
premisas sean suficientemente grandes como para que el cociente de Jaccard sea
estable; construirlos es trabajo futuro (sección~\ref{sec:conclusion}).

\textbf{D3 mide formalizaciones, no ideas.} El caso T07 expone un límite de la métrica que excede la elección del umbral. Mathlib contiene típicamente una sola prueba canónica de cada teorema clásico, de modo que dos estrategias que un matemático consideraría distintas pueden converger al mismo lema al ser formalizadas. La consecuencia es que D3 mide distancia entre \textit{formalizaciones}, y la relación entre esa distancia y la distancia entre las ideas informales subyacentes queda sin establecer: una distancia de 1{,}0 puede deberse a que el formalizador eligió lemas distintos, y una distancia baja a que convergió al mismo lema de Mathlib. Reportamos D3 como relativa a la formalización sobre Mathlib v4.29.0. Cuantificar con qué frecuencia ocurre esa convergencia requiere un experimento sistemático sobre múltiples teoremas y múltiples formalizadores, que queda como trabajo futuro.

\subsection{D2: filtro de trivialidad}

D2 se evaluó sobre 24 de los 26 teoremas del conjunto de evaluación. Dos casos (T20 y T21) no fueron formalizados en esta corrida por razones ajenas al instrumento. Los 24 evaluados cubren triviales por diseño, clásicos presentes en Mathlib, pares con distinta prueba, enunciados cercanos, casos generados por un LLM y casos de falla deliberados.

\begin{table}[htbp]
\centering
\small
\setlength{\tabcolsep}{5pt}
\begin{tabular}{llll}
\toprule
Caso & Enunciado & Táctica & Expectativa \\
\midrule
T14 & suma de cuatro pares es par        & \texttt{aesop}    & trivial \\
T15 & $2 + 2 = 4$                        & \texttt{decide}   & trivial \\
T16 & $n + 0 = n$                        & \texttt{norm\_num} & trivial \\
T17 & $n \leq n + 1$                     & \texttt{norm\_num} & trivial \\
T19 & generado por LLM sobre pares       & \texttt{aesop}    & ambigua \\
T22 & $n$ par $\Rightarrow n+0$ par      & \texttt{norm\_num} & diseñado para D1 \\
\bottomrule
\end{tabular}
\caption{Los seis enunciados que D2 cerró con una táctica de \texttt{T\_AUTO}. Los cuatro primeros tenían expectativa de trivialidad definida; T19 y T22 se discuten en el texto.}
\label{tab:d2trivial}
\end{table}

T19 fue generado por un LLM con la instrucción de enunciar y demostrar un teorema original sobre números pares; D2 lo cerró con \texttt{aesop}, lo que es consistente con la hipótesis de que los modelos tienden a producir enunciados triviales, pero su expectativa en el conjunto de evaluación no era binaria. T22 es lógicamente equivalente a ``si $n$ es par entonces $n$ es par''; D2 lo cerró correctamente, pero el caso fue construido para poner a prueba la equivalencia sintáctica en D1, no la trivialidad. Ninguno de los dos admite una clasificación de acierto o error para D2.

\begin{table}[htbp]
\centering
\small
\setlength{\tabcolsep}{5pt}
\begin{tabular}{p{4.6cm}p{6.4cm}l}
\toprule
Casos & Categoría & Resultado D2 \\
\midrule
T01--T06                & clásicos presentes en Mathlib            & no trivial \\
T07a, T08a, T09a        & pares con distinta demostración          & no trivial \\
T10--T13                & enunciados cercanos (primos impares, AM--GM) & no trivial \\
T18                     & suma de los $n$ primeros impares $= n^2$ & no trivial \\
T23                     & grafo conexo y acíclico es árbol         & no trivial \\
T24                     & haces coherentes sobre esquema noetheriano & no trivial \\
T25                     & $n$ par $\iff 2 \mid n$                  & no trivial \\
T26                     & suma de $n$ pares es par                 & no trivial \\
\bottomrule
\end{tabular}
\caption{Los 18 enunciados que D2 clasificó correctamente como no triviales.}
\label{tab:d2notrivial}
\end{table}

\textbf{Casos notables.} T18 es una trampa de control: requiere inducción y D2 correctamente no lo cerró. T23 tampoco fue cerrado en esta corrida, a diferencia de corridas anteriores en que \texttt{tauto} lo cerraba sobre una definición conjuntiva; el enunciado Lean efectivamente usado aquí difiere del que producía ese falso positivo. T14, en cambio, sí fue cerrado por \texttt{aesop} en esta corrida, mientras que en corridas previas con un enunciado más complejo la misma táctica excedía el presupuesto por más de un orden de magnitud. Ambos casos ilustran la sensibilidad de D2 al enunciado concreto que produce el formalizador, discutida en la sección~\ref{sec:limitaciones}.

\subsubsection*{Agregado}

Sobre los 24 teoremas evaluados, D2 acierta en 22 (91{,}7\,\%). T19 y T22 permanecen en el denominador pero no en el numerador, por carecer de expectativa binaria contra la cual medir acierto; sobre los 22 casos con expectativa definida no se registran falsos positivos ni falsos negativos.

\subsection{D1 C\_F: cobertura del corpus formal}

La rama C\_F encontró match para los 18 teoremas no triviales que alcanzaron esa etapa. Los seis restantes no registran match porque D1 nunca se ejecutó sobre ellos: son exactamente los seis que D2 detectó como triviales, sobre los cuales el árbol aplicó cortocircuito. La ausencia de match en esos casos refleja no-ejecución, no un fallo de la búsqueda.

\textbf{Este resultado no debe leerse como cobertura medida.} El criterio de match de C\_F acepta el primer resultado que devuelve Leandex sin aplicar umbral de similitud, porque la versión actual de la API no proporciona scores (sección~\ref{sec:pipeline}). Un criterio así solo puede devolver ausencia de match cuando la API responde vacío, de modo que 18 matches sobre 18 consultas no distingue cobertura real de ausencia de filtro. Los 18 casos corresponden a teoremas clásicos que efectivamente están en Mathlib, y la composición del conjunto de evaluación —dominada por resultados canónicos— hace esperable ese resultado con cualquier criterio. Establecer la precisión de C\_F requiere un umbral de similitud o una verificación manual de cada match, y ninguna de las dos está disponible en esta versión.

La rama C\_I no se activó para ningún teorema. La causa no es el umbral de embeddings sino la condición de entrada: C\_I solo se consulta cuando D2 es negativo y C\_F no encontró match, y esa combinación nunca se produjo. Todos los no triviales tuvieron match en C\_F, y todos los que no lo tuvieron fueron detenidos en D2. El conjunto de evaluación, por composición, no ejercita esa rama.

\subsection{Resumen}

\begin{table}[htbp]
\centering
\small
\setlength{\tabcolsep}{4pt}
\renewcommand{\arraystretch}{1.15}
\begin{tabular}{p{1.9cm}p{2.6cm}p{6.4cm}p{3.4cm}}
\toprule
Instrumento & Ground truth & Resultado & Estado \\
\midrule
D3 &
5 pares con etiqueta humana previa &
5/5 consistentes con el juez bajo $\theta = 0{,}5$. Solo T08 y T07 restringen el
umbral, y T07 es degenerado ($|\cup| = 2$): la evidencia sostiene $\theta < 0{,}7222$, no un valor puntual. &
Funcional; $\theta$ no calibrado \\
\addlinespace
D2 &
24 teoremas del conjunto de evaluación &
22/24 aciertos (91{,}7\,\%); T19 y T22 en el denominador sin expectativa binaria. Sin falsos positivos ni negativos sobre los 22 restantes. &
Funcional \\
\addlinespace
D1 C\_F &
18 teoremas no triviales &
18 matches sobre 18 consultas, pero sin umbral de similitud el resultado no distingue cobertura de ausencia de filtro. &
No concluyente \\
\addlinespace
D1 C\_I &
--- &
Rama no alcanzada: la condición de entrada nunca se cumplió con este conjunto. &
No evaluada \\
\bottomrule
\end{tabular}
\caption{Estado de validación de cada instrumento. Solo D2 admite una medida de acierto interpretable con el conjunto de evaluación actual.}
\label{tab:instrumentos}
\end{table}


\section{Experimentos}
\label{sec:experimentos}

Esta sección reporta la evaluación del pipeline sobre artículos reales. El ground truth proviene de artículos retirados de arXiv cuyo comentario de retiro declara duplicación de resultados previos. La evaluación produjo tanto medidas de funcionamiento como la caracterización de tres obstáculos que limitan el enfoque con independencia de esta implementación.

\subsection{Construcción del conjunto de artículos retirados}

Se consultó la API pública de arXiv con una búsqueda de artículos retirados en categorías de matemática, que devolvió aproximadamente 2.600 resultados. Sobre el texto del comentario de retiro se aplicaron 23 patrones de expresión regular diseñados para identificar retiros por duplicación de resultados previos y excluir retiros por errores, lagunas o problemas administrativos. Los patrones más frecuentes corresponden a fórmulas como \textit{result was already known}, \textit{had already been proved by}, \textit{results are not new} y \textit{subsumed by}.

El filtrado produjo 33 candidatos, de los cuales 26 resultaron viables en el momento de la construcción: fuente LaTeX descargable y al menos un entorno de teorema detectable. Los 7 restantes usan formatos que el parser no reconoce, como AMS-\TeX{} o nombres abreviados de entorno. Los 26 viables cubren 17 categorías, con concentración en combinatoria y geometría algebraica, y un rango de años de 2001 a 2026. Para cada uno se seleccionaron dos controles de la misma categoría, con año de publicación dentro de $\pm 1$, no retirados y con fuente verificable; se examinaron 382 candidatos para obtener 52 controles.

\textbf{Calidad del ground truth.} De los 26 artículos retirados, 12 identifican explícitamente el trabajo previo que duplica su resultado; los 14 restantes usan fórmulas genéricas sin especificar la fuente. Entre los 12 que sí lo identifican, solo una minoría proporciona un identificador de arXiv: la mayoría cita por autor y publicación, lo que exige resolución manual. Esa asimetría tiene consecuencias que se analizan en la sección~\ref{ssec:cobertura}.

\subsection{Disponibilidad de las fuentes: un obstáculo estructural}
\label{ssec:fuentes}

Durante la ampliación del conjunto se descubrió una restricción que afecta al método completo: \textbf{arXiv elimina el código fuente de un artículo al retirarlo.} El endpoint de fuente devuelve 404 para artículos retirados, sin excepciones observadas.

Se verificó sobre una muestra aleatoria de 25 artículos retirados por duplicación extraídos de un conjunto mayor de candidatos: ninguno conservaba fuente descargable. La hipótesis de que los artículos marcados como subsumidos pero no retirados formalmente conservarían su fuente resultó falsa: en arXiv, la marca de subsunción implica retiro.

Las consecuencias son dos. Primero, ningún pipeline que parta del fuente LaTeX puede procesar artículos retirados, salvo que las fuentes se hayan descargado antes del retiro. Segundo, y más grave para la evaluación, el conjunto que construimos no es reproducible desde arXiv: los artículos que eran viables en el momento de la construcción dejaron de serlo, y su material procesable sobrevive únicamente en copias locales. Cualquier trabajo futuro que use artículos retirados como ground truth de duplicación debe contemplar esta restricción desde el diseño, conservando las fuentes en el momento de la selección.

\subsection{Selección del modelo de formalización}

Cuatro modelos de lenguaje se compararon en la tarea de formalización \textit{statement-only} sobre cinco artículos retirados.

\begin{table}[htbp]
\centering
\small
\begin{tabular}{lcl}
\toprule
Modelo & Éxito & Modo de falla predominante \\
\midrule
Qwen 3.7-max      & 5/5 & --- \\
GLM-5.2           & 3/5 & errores reparables, un timeout \\
DeepSeek V4 Pro   & 0/5 & errores de síntesis, definición placeholder \\
DeepSeek V4 Flash & 0/5 & errores de síntesis, tokens inesperados \\
\bottomrule
\end{tabular}
\caption{Comparación de backends de formalización sobre cinco enunciados. Qwen 3.7-max fue el único que produjo definiciones matemáticamente sustantivas en los cinco casos, y es el modelo usado en el resto de los experimentos.}
\label{tab:modelos}
\end{table}

\textbf{Advertencia sobre esta selección.} Los cinco artículos del banco de comparación son los mismos que integran el grupo de retirados del estudio profundo. El modelo fue elegido, por lo tanto, sobre parte del conjunto de evaluación, lo que introduce un sesgo optimista sobre su desempeño en ese grupo y ofrece una explicación alternativa a la asimetría de fidelidad reportada más abajo.

\subsection{Estudio profundo}

El estudio profundo evaluó 10 artículos —5 retirados y 5 controles emparejados— con formalización \textit{statement-only} del enunciado, filtro de trivialidad, y búsqueda en corpus formal e informal. Los retirados se seleccionaron manualmente entre los 26 viables como casos claros de duplicación.

De los 10 artículos, 7 se formalizaron con éxito. De los 3 restantes, uno falló por error de compilación y dos por exceder el límite del modelo con enunciados de cuatro niveles de casos anidados.

\subsubsection{Auditoría de fidelidad}

El primer autor revisó manualmente las 7 formalizaciones exitosas, comparando el código Lean generado contra el enunciado original.

\begin{table}[htbp]
\centering
\small
\begin{tabular}{llll}
\toprule
Artículo & Rol & Fidelidad & Problema detectado \\
\midrule
1609.02090v1   & retirado & fiel & --- \\
1207.0631v1    & retirado & fiel & --- \\
1212.0196v1    & retirado & fiel & --- \\
1004.3381v1    & retirado & fiel$^{\dagger}$ & --- \\
1101.3720v1    & control  & incorrecta & \texttt{sorry} en la definición central \\
0904.1783v3    & control  & aproximación & solo una dirección de una equivalencia \\
math/0504586v2 & control  & incorrecta & evento y medida trivializados \\
\bottomrule
\end{tabular}
\caption{Auditoría manual de las siete formalizaciones exitosas. $^{\dagger}$Ver nota en el texto sobre este caso.}
\label{tab:fidelidad}
\end{table}

La tasa de fidelidad es de 4/4 entre los retirados y 0/3 entre los controles. Dos explicaciones compiten. La primera es un sesgo de complejidad: los controles, emparejados por categoría y año pero no por dificultad de enunciado, tienden a enunciados más largos y con más casos anidados, hipótesis que refuerzan los dos controles que fallaron por exceder el límite del modelo. La segunda, más simple, es el sesgo de selección del modelo señalado arriba: el formalizador fue elegido por su desempeño sobre estos mismos retirados y sobre ningún control. Los datos disponibles no permiten distinguir entre ambas.

La consecuencia experimental es que los tres controles evaluados no aportan evidencia sobre falsos positivos: sus formalizaciones no representan los teoremas originales, de modo que cualquier veredicto sobre ellos es informativo sobre la formalización y no sobre el detector.

\subsubsection{Cinco modos de falla de la verificación por compilación}
\label{ssec:modosfalla}

La formalización automática introduce un problema que no existe en la verificación manual: un archivo Lean puede compilar sin errores sin contener una formalización fiel del teorema original. Durante el desarrollo se documentaron cinco modos de falla sucesivos. Cada uno motivó una guardia nueva, y cada guardia dejó pasar el siguiente.

\begin{enumerate}
\item \textbf{Archivo vacío.} Un archivo sin declaraciones compila limpiamente. En la primera corrida, cinco de cinco formalizaciones eran archivos vacíos que satisfacían el criterio de éxito original. La guardia correspondiente exige la presencia de al menos una declaración sustantiva.

\item \textbf{Definición placeholder.} Una definición de la forma \texttt{def CongruentNumber := True} compila, contiene una declaración sustantiva y satisface la guardia anterior, pero no captura ninguna definición matemática. La respuesta fue reforzar el prompt contra la trivialización de definiciones.

\item \textbf{Definición correcta sin el teorema.} Tras ese refuerzo, el mismo artículo produjo una definición matemáticamente correcta pero omitió el enunciado del teorema. El archivo compilaba, la definición era fiel, y el objeto a evaluar estaba ausente.

\item \textbf{\texttt{sorry} en una definición auxiliar.} Un \texttt{sorry} dentro del cuerpo de una definición compila y satisface cualquier guardia sintáctica sobre las declaraciones: la firma existe, el contenido es un marcador. El criterio de éxito rechaza archivos con \texttt{sorry}, pero solo si el verificador lo detecta en el lugar correcto.

\item \textbf{Una sola dirección de una equivalencia.} Un teorema enunciado como equivalencia en el artículo original se formalizó en una sola dirección. El archivo compila, las definiciones son fieles, la dirección formalizada es correcta, y falta la mitad del enunciado.
\end{enumerate}

Los cinco peldaños comparten una propiedad: la compilación verifica coherencia, nunca fidelidad semántica. Cada guardia cierra una clase de fallo y deja pasar la siguiente porque la pregunta de fondo —si el código Lean generado significa lo mismo que el enunciado original— no es decidible por medios sintácticos. De ahí que la auditoría humana sea irreducible: ningún test automático puede establecer que definir un evento de percolación como el conjunto total no es una formalización aceptable.

Esta observación tiene alcance más allá de AViD. Los pipelines de formalización automática reportan habitualmente tasas de compilación exitosa como medida de desempeño. Los cinco modos de falla muestran que esa métrica y la fidelidad semántica son magnitudes distintas, y que la primera no acota a la segunda. Qué fracción de las formalizaciones que compilan es semánticamente infiel es una pregunta abierta que requiere auditoría sistemática con protocolo ciego, y que este trabajo no responde.

\subsection{Detección de duplicación}

\subsubsection{Dos defectos corregidos}

Dos defectos del pipeline se identificaron ejecutando la evaluación, y ambos producían falsos positivos de duplicación.

\textbf{Ausencia de orden temporal.} El artículo 1207.0631v1, publicado en 2012, recibió veredicto de \texttt{CONOCIDO\_LITERATURA} sobre la base de un candidato de 2018 que el juez clasificó como equivalente. Un trabajo posterior no puede ser aquello que un trabajo anterior duplica: el veredicto era lógicamente imposible. La incorporación del filtro temporal (sección~\ref{ssec:d1}) descarta los tres candidatos posteriores a la fecha del artículo y el veredicto pasa correctamente a \texttt{NOVEDAD\_ENUNCIADO}.

\textbf{Recuperación del propio artículo.} Los índices consultados contienen los artículos evaluados. Sin exclusión explícita, tres de los siete artículos se recuperaron a sí mismos como candidatos y el juez los clasificó como equivalentes o especializaciones, produciendo veredictos de no-novedad sin contenido. Con exclusión, los tres veredictos se corrigen y no se observan autocoincidencias.

Ambos defectos son genéricos: afectan a cualquier sistema que verifique novedad por recuperación sobre índices que contienen el material evaluado. El segundo, además, había sido corregido en una rama del sistema y permanecía activo en otra, lo que ilustra que una corrección puntual no se propaga sola.

\subsubsection{Resultados}

Con ambas correcciones y consultas construidas a partir del enunciado, el pipeline detectó una coincidencia legítima —un artículo de control cuyo resultado aparece en un trabajo anterior efectivamente indexado— y ninguna sobre los artículos retirados con formalización fiel.

% [PENDIENTE] Confirmar tras la auditoría del match 0904.1783 vs 0405089:
% si es legítimo, es evidencia de funcionamiento; si el juez se equivocó,
% es el primer falso positivo documentado y debe reportarse como tal.
% Ajustar el párrafo anterior segun el resultado.

La ausencia de detecciones entre los retirados no es un fallo del detector. Los duplicadores de esos artículos —resultados de Hardy y Littlewood, Monsky, y Gyárfás y Lehel— son anteriores a la era de los preprints electrónicos. La sección~\ref{ssec:cobertura} muestra que ese hecho, y no la métrica ni el juez, determina el resultado.

\subsubsection{Similitud e identidad de resultado}

Un caso ilustra la distancia entre recuperar un trabajo similar y recuperar el trabajo duplicado. El artículo 1609.02090v1 contiene dos resultados independientes: uno que el propio artículo atribuye explícitamente a un trabajo previo, y otro —el motivo del retiro— que había sido demostrado por Hardy y Littlewood alrededor de 1928. La búsqueda recuperó el primero, es decir, el trabajo que duplica un resultado que el autor nunca reclamó como propio, y no el segundo. La coincidencia es correcta y a la vez irrelevante para la pregunta formulada.

\subsection{Cobertura de los corpus de recuperación}
\label{ssec:cobertura}

Los resultados anteriores sugieren que el factor limitante no está en el pipeline sino en qué contienen los índices. Para verificarlo se midió la cobertura directamente sobre los duplicadores conocidos, sin intervención del pipeline: la pregunta es si el trabajo duplicado está indexado, no si el sistema lo encuentra.

\textbf{Recuperabilidad de los duplicadores.} De los duplicadores identificados en los comentarios de retiro, la mayoría corresponde a trabajos anteriores a 1991: resultados clásicos publicados en revistas, no preprints. Solo una minoría tiene identificador de arXiv. Esta distribución no es accidental: cuando un resultado se re-demuestra sin advertirlo, el trabajo original suele ser un resultado establecido y antiguo, precisamente el tipo de material que los índices de teoremas cubren peor.

% [PENDIENTE — verificacion del metodo de consulta]
% Insertar aqui las cifras de cobertura una vez confirmado si la busqueda
% por titulo en un indice de ENUNCIADOS produce falsos negativos. Si el
% resultado 0/N sobre duplicadores post-1991 se sostiene con busqueda por
% enunciado, es el hallazgo principal de esta subseccion y debe reportarse
% con las cifras. Si es artefacto del metodo, remedir antes de afirmar nada
% sobre la cobertura de indices de terceros.

\textbf{Consecuencia.} Un duplicador ausente del índice impone un techo de cero al recall, con independencia de la calidad de la métrica de similitud, del umbral de embeddings o del juez. Ninguna mejora del pipeline modifica ese techo. La vía de mejora es la ampliación del corpus, y la incorporación de un índice de revistas revisadas por pares con cobertura desde 1826 no fue suficiente para los casos evaluados. Esta limitación no es específica de AViD: cualquier sistema que verifique novedad por recuperación la enfrenta con la infraestructura actualmente disponible.

\subsection{Un baseline de similitud semántica}

Como referencia, se evaluó si una métrica de similitud semántica aplicada directamente al texto de los enunciados, sin formalización ni árbol de decisión, permite distinguir artículos retirados de controles. Sobre 37 artículos evaluables, una prueba de Mann-Whitney sobre las distribuciones de puntajes no arroja diferencia significativa ($p = 0{,}854$).

Este resultado no debe leerse como evidencia de ausencia de efecto. El tamaño de muestra es pequeño y desigual, la exclusión de artículos no fue aleatoria —los controles tuvieron aproximadamente el doble de probabilidad de quedar fuera por ausencia de enunciado extraíble—, y no se reporta poder estadístico. Lo que el resultado sí indica es que la similitud textual entre enunciados, tomada aisladamente, no es un discriminante útil de duplicación en este conjunto.

Una versión previa de este mismo baseline arrojó un resultado aparentemente positivo que resultó artefactual: sin exclusión del propio artículo, la mayoría de las coincidencias fuertes eran autocoincidencias. Se documenta aquí como advertencia para trabajos que evalúen sobre índices que contienen el material evaluado.

\subsection{Resumen}

\begin{table}[htbp]
\centering
\small
\setlength{\tabcolsep}{4pt}
\renewcommand{\arraystretch}{1.15}
\begin{tabular}{p{3.6cm}p{1.6cm}p{8.6cm}}
\toprule
Evaluación & $n$ & Resultado \\
\midrule
Selección de backend & 5 $\times$ 4 & Qwen 3.7-max 5/5; los dos modelos DeepSeek 0/5. Selección contaminada por solapamiento con el conjunto de evaluación. \\
\addlinespace
Formalización & 10 & 7 formalizados. Fidelidad 4/4 en retirados, 0/3 en controles; los controles no aportan evidencia sobre falsos positivos. \\
\addlinespace
Modos de falla & --- & Cinco modos sucesivos en que la compilación exitosa no implica fidelidad semántica. \\
\addlinespace
Detección de duplicación & 7 & Dos defectos corregidos (orden temporal, autocoincidencia). Sin detecciones sobre retirados con duplicador anterior a 1991. \\
\addlinespace
Cobertura de corpus & --- & El contenido de los índices, y no la métrica, determina el techo de recall. \\
\addlinespace
Disponibilidad de fuentes & 25 & Ninguna fuente de artículo retirado sigue disponible en arXiv. \\
\addlinespace
Baseline semántico & 37 & Sin diferencia significativa ($p = 0{,}854$); muestra pequeña y sesgada. \\
\bottomrule
\end{tabular}
\caption{Resumen de la evaluación. Las tres últimas filas corresponden a obstáculos estructurales, no a mediciones de desempeño del pipeline.}
\label{tab:resumen}
\end{table}


\section{Conclusión}
\label{sec:conclusion}

\subsection{Lo construido}

AViD Journal recorre el camino completo desde un artículo en LaTeX hasta un veredicto de novedad: parser de entornos matemáticos con ordenamiento topológico de dependencias, formalización en Lean 4 con abstracción sobre ocho backends, tres corpus de referencia —Mathlib vía Leandex, TheoremSearch y Matlas—, y un árbol de decisión de ocho veredictos sobre tres dimensiones evaluadas en orden de costo creciente. El sistema corre sobre Windows nativo, incluida la extracción de premisas, sin requerir WSL ni Docker.

\subsection{Lo medido}

D2 acertó en 22 de 24 teoremas del conjunto de evaluación (91{,}7\,\%), sin falsos positivos ni falsos negativos sobre los casos con expectativa definida. D3 resultó consistente con el juicio humano en los cinco pares de calibración, aunque solo dos de ellos restringen el umbral. D1 C\_F devolvió match en las 18 consultas que alcanzaron esa etapa, pero sin umbral de similitud ese resultado no distingue cobertura real de ausencia de filtro. Sobre artículos reales, el pipeline detecta duplicación cuando el trabajo duplicado está indexado en el corpus, y no la detecta cuando no lo está.

La comparación de cuatro modelos de formalización cuantifica la viabilidad relativa de distintos backends y motivó la elección del que se usa en los experimentos. Se documenta además un defecto metodológico que afecta a cualquier evaluación sobre índices de teoremas: sin exclusión explícita del propio artículo, un paper se recupera a sí mismo como candidato y el juez lo clasifica como equivalente, produciendo falsos positivos de duplicación.

\subsection{Lo aprendido}

La evaluación produjo un resultado más informativo que cualquier medida de desempeño del pipeline: la identificación de tres obstáculos que no son deficiencias de esta implementación sino propiedades del problema con la infraestructura hoy disponible.

El primero es que la compilación verifica coherencia, nunca fidelidad semántica: un archivo Lean puede compilar sin errores y no representar el enunciado original, y ninguna guardia sintáctica cierra esa brecha (sección~\ref{sec:experimentos}). El segundo es que la cobertura de los índices de teoremas, y no la métrica de similitud ni el juez, impone el techo de recuperación: un duplicador ausente del corpus no puede ser encontrado por ningún pipeline, por bueno que sea. El tercero es que arXiv elimina el código fuente de los artículos al retirarlos, lo que restringe qué puede procesarse y compromete la reproducibilidad de cualquier benchmark construido sobre ellos: el conjunto de artículos que reunimos es replicable en su composición —identificadores y comentarios de retiro— pero no en su material procesable, que debe conservarse localmente antes del retiro.

Los tres son transferibles: cualquier sistema que verifique novedad por recuperación y formalización va a encontrarlos.

\subsection{Limitaciones y trabajo futuro}

Las limitaciones restantes de esta versión definen, en su mayoría, el trabajo que sigue.

\textbf{Alcance de la métrica.} AViD evalúa cada enunciado por separado, de modo que un artículo cuya novedad resida en la combinación de resultados conocidos no sería detectado; una capa de agregación es trabajo futuro. D3 opera solo sobre \texttt{Prop} y asume irrelevancia de pruebas: extenderla a \texttt{Type}, donde la noción adecuada sería homotópica, es un objetivo de largo plazo. Más de fondo, el marco mide novedad sobre enunciados y pruebas congelados en un corpus y no captura la novedad conceptual de una definición que reformula un problema o de un método transferible a otros dominios~\cite{lakatos1976}; y la identidad de pruebas es indecidible en general~\cite{dosen2003}, por lo que Jaccard sobre premisas es una aproximación de ingeniería, con falsos positivos cuando pruebas equivalentes usan premisas superficialmente distintas y falsos negativos cuando estrategias distintas comparten lemas del núcleo.

\textbf{Calibración de D3.} El umbral $\theta = 0{,}5$ es una elección de diseño. La escalera de calibración es consistente con el juicio humano en los cinco pares, pero solo dos restringen el umbral y uno de ellos es degenerado, de modo que la evidencia acota el intervalo admisible a $\theta < 0{,}7222$ sin fijar un valor. Calibrar requiere un corpus de pares de demostraciones del mismo teorema con etiqueta humana previa, seleccionados de modo que sus distancias caigan en el interior del rango y sus conjuntos de premisas sean lo bastante grandes como para que el cociente sea estable; un corpus así permitiría además ponderar las premisas por IDF~\cite{huch2022} en lugar de contarlas por igual, sustituir los conjuntos de premisas por grafos de dependencia completos, y reportar precisión y exhaustividad de D3 sobre el eje de identidad de pruebas. Esto es lo primero que haremos en la próxima versión.

\textbf{Dependencia del formalizador.} AViD evalúa la demostración formalizada, no la informal. Una distancia máxima puede deberse a que el formalizador eligió lemas distintos, y una distancia baja a que convergió al mismo lema de Mathlib (sección~\ref{ssec:d3cal}). El mismo efecto opera sobre D2, cuyo resultado depende del par (enunciado formalizado, táctica, presupuesto) y no del teorema. Cuantificar con qué frecuencia la autoformalización preserva o colapsa la distinción entre estrategias requiere un experimento sistemático sobre múltiples teoremas y múltiples modelos, con auditoría humana de la fidelidad de cada formalización; es la línea de trabajo que consideramos más prometedora. Ejecutar varios formalizadores en paralelo reduciría además la dependencia de un único proveedor.

\textbf{Precisión de D1.} La comparación de enunciados es sintáctica: incorporar \texttt{isDefEq} permitiría reconocer equivalencias definicionales que hoy producen falsos negativos. La rama informal depende además de poder formalizar demostraciones ajenas, capacidad que tuvo 0\,\% de éxito en la prueba de concepto; sin ese eslabón D3 no puede aplicarse a matches del corpus informal.
Dos extensiones apuntan a escenarios de publicación acelerada como el descrito
en la introducción. La primera es de granularidad temporal: la fecha que el
filtro deriva del identificador de arXiv resuelve a nivel de mes, suficiente
para descartar candidatos posteriores por años pero no para establecer
prioridad entre trabajos aparecidos con días de diferencia; sustituirla por la
fecha de envío que provee la API resolvería ese caso. La segunda es de alcance:
el pipeline formaliza enunciados y compila el resultado, de modo que la
infraestructura para evaluar la validez de un contraejemplo ---una tarea
computacionalmente mucho más barata que verificar una demostración--- ya está
presente, aunque el árbol de decisión actual solo emite veredictos de novedad.
Extenderlo en esa dirección requeriría medir primero con qué fidelidad el
formalizador traduce enunciados de contraejemplo, y por lo tanto enfrenta el
mismo obstáculo documentado en la sección~\ref{ssec:modosfalla}.

\textbf{Escala de la evaluación.} Los 24 teoremas del conjunto de evaluación fueron construidos por el primer autor, y el estudio profundo operó sobre artículos seleccionados entre los que resultaron parseables y formalizables, lo que introduce un sesgo de accesibilidad técnica cuyo efecto sobre la duplicabilidad desconocemos. Tres direcciones amplían esa escala: evaluar D2 contra un corpus de conjeturas etiquetadas de forma independiente, lo que daría una medida con ground truth externo; ejecutar el pipeline completo sobre Mathlib, que ofrece un corpus formal de magnitud muy superior al conjunto actual; y construir un benchmark sobre artículos de arXiv autoformalizados, donde la fuente sigue disponible. En cuanto al ground truth de duplicación, el comentario de retiro es un proxy imperfecto: de los artículos retirados viables, menos de la mitad identifica explícitamente el trabajo que los duplica, y entre los que lo hacen la mayoría lo cita por autor y publicación en lugar de por identificador, lo que exige una resolución manual con su propio margen de error.

\textbf{Ingeniería.} La extracción de premisas requiere cargar Mathlib completo, de modo que D3 se ofrece como análisis a pedido y no dentro del flujo interactivo; los imports de módulos específicos no funcionan sobre archivos temporales, lo que obliga a cargar la biblioteca entera en cada invocación de D2.

\subsection{Cierre}

AViD Journal no resuelve el problema de la verificación automática de novedad. Lo plantea de manera operacional, lo implementa de extremo a extremo, y documenta con precisión dónde falla. El sistema opera sobre el subconjunto de teoremas expresables en el fragmento \texttt{Prop} de Lean 4, con enunciado parseable y formalizable por alguno de los modelos disponibles; dentro de ese perímetro el árbol de decisión produce veredictos verificables, y fuera de él es explícitamente incompleto.

La tesis del trabajo es que la pregunta sobre si puede arbitrarse novedad automáticamente no se responde con un diseño de métrica ni con un argumento de arquitectura, sino construyendo el sistema, ejecutándolo contra ground truth real, y reportando qué funcionó y qué no. Los tres obstáculos identificados sugieren que el cuello de botella no está donde el diseño de la métrica invita a buscarlo.


\begin{thebibliography}{99}

\bibitem{abouzaid2026}
M.\textasciitilde{}Abouzaid \emph{et al.}, ``First Proof,''
arXiv:2602.05192, 2026.

\bibitem{firstproof2b}
M.\textasciitilde{}Abouzaid, N.\textasciitilde{}Srivastava, R.\textasciitilde{}Ward, and L.\textasciitilde{}Williams,
``First Proof: Second Batch,'' First Proof Foundation, junio 2026.
\url{https://1stproof.org/assets/docs/report.pdf}

\bibitem{tao2025}
T.\textasciitilde{}Tao, ``Machine-Assisted Proof,''
\emph{Notices of the AMS}, 2025.

\bibitem{kasaura2025}
K.\textasciitilde{}Kasaura \emph{et al.}, ``Discovering New Theorems via LLMs with In-Context Proof Learning in Lean,''
arXiv:2509.14274, 2025.

\bibitem{theoremsearch2026}
A.\textasciitilde{}Ilin, T.\textasciitilde{}Alper, and L.\textasciitilde{}Inchiostro, ``TheoremSearch: Semantic Search over Theorems,''
arXiv:2602.05216, 2026.

\bibitem{theoremgraph2026}
TheoremGraph + LeanGraph, arXiv:2606.25363, 2026. {[}VERIFICAR{]}

\bibitem{compose2026}
COMPOSE, arXiv:2605.30333, 2026. {[}VERIFICAR{]}

\bibitem{leanconjecturer2025}
N.\textasciitilde{}Onda \emph{et al.}, ``LeanConjecturer: Filtering Novelty in LLM-Generated Conjectures,''
arXiv:2506.22005, 2025.

\bibitem{matlas2026}
Matlas: 8.07M theorem statements from 435K peer-reviewed articles (1826--2025),
arXiv:2604.17484, 2026.

\bibitem{yoo2025}
S.\textasciitilde{}Yoo, ``The Axiom-Based Atlas of Mathematical Theorems,''
arXiv:2504.00063, 2025.

\bibitem{kaliszyk2013}
C.\textasciitilde{}Kaliszyk and J.\textasciitilde{}Urban, ``MaSh: Machine Learning for Sledgehammer,''
\emph{Interactive Theorem Proving (ITP)}, 2013. {[}VERIFICAR{]}

\bibitem{magnushammer2024}
M.\textasciitilde{}Mikula, A.\textasciitilde{}Jiang, W.\textasciitilde{}Li \emph{et al.}, ``Magnushammer: A Transformer-Based Approach to Premise Selection,''
\emph{ICLR}, 2024.

\bibitem{piotrowski2023}
B.\textasciitilde{}Piotrowski \emph{et al.}, ``Machine-Learned Premise Selection for Lean,''
arXiv:2304.00994, 2023.

\bibitem{network2026}
``The Network Structure of Mathlib,''
arXiv:2604.24797, 2026. {[}VERIFICAR{]}

\bibitem{huch2022}
F.\textasciitilde{}Huch, ``Structure in Theorem Proving: Analyzing the Archive of Formal Proofs,''
arXiv:2209.13305, 2022.

\bibitem{mining2024}
``Mining Math Conjectures from LLMs: A Pruning Approach,''
arXiv:2412.16177, 2024.

\bibitem{synthetic2024}
``Synthetic Theorem Generation in Lean,''
OpenReview EeDSMy5Ruj, 2024.

\bibitem{withdrarxiv2024}
R.\textasciitilde{}Rao \emph{et al.}, ``WithdrarXiv: A Large-Scale Dataset of Retracted Papers,''
arXiv:2412.03775, 2024.

\bibitem{pseudo2026}
Pseudo-Formalization / ArxivMathGradingBench, 2026. {[}VERIFICAR{]}

\bibitem{merlean2026}
MerLean, 2026. {[}VERIFICAR{]}

\bibitem{survey2026}
``AI for Mathematics: Progress, Challenges, and Prospects,''
arXiv:2601.13209, 2026.

\bibitem{lakatos1976}
I.\textasciitilde{}Lakatos, \emph{Proofs and Refutations}, Cambridge University Press, 1976.

\bibitem{dosen2003}
K.\textasciitilde{}Dosen, ``Identity of Proofs,''
\emph{Bulletin of Symbolic Logic}, vol.\textasciitilde{}9, no.\textasciitilde{}4, 2003.

\bibitem{theoremsearch2026}
L.~Alexander, E.~Leonen, S.~Szeto, A.~Remizov, I.~Tejeda, J.~Alper,
G.~Inchiostro, and V.~Ilin, ``Semantic Search over 9 Million Mathematical Theorems,''
arXiv:2602.05216, 2026.

\bibitem{matlas2026}
``Matlas: A Semantic Search Engine for Mathematics,'' arXiv:2604.17484, 2026.

\bibitem{network2026}
X.~Li, N.~Peng, S.~Severini, and P.~Shafto, ``The Network Structure of Mathlib,''
arXiv:2604.24797, 2026.

\bibitem{leanexplore2025}
J.~Asher, ``LeanExplore: A Search Engine for Lean 4 Declarations,''
arXiv:2506.11085, 2025.

\bibitem{leanfinder2025}
J.~Lu, K.~Emond, W.~Sun, and W.~Chen, ``Lean Finder: Semantic Search for Mathlib
That Understands User Intents,'' 2nd AI for Math Workshop @ ICML, 2025.

\bibitem{kuhlwein2013}
D.~K\"uhlwein, J.~C.~Blanchette, C.~Kaliszyk, and J.~Urban,
``MaSh: Machine Learning for Sledgehammer,''
\emph{Interactive Theorem Proving (ITP 2013)}, LNCS 7998, pp.~35--50, Springer, 2013.

\bibitem{kaliszyk2014}
C.~Kaliszyk and J.~Urban, ``Learning-Assisted Automated Reasoning with Flyspeck,''
\emph{Journal of Automated Reasoning}, vol.~53, no.~2, pp.~173--213, 2014.

\bibitem{magnushammer2024}
M.~Miku{\l}a et al., ``Magnushammer: A Transformer-Based Approach to Premise Selection,''
arXiv:2303.04488, \emph{ICLR}, 2024.

\end{thebibliography}

\end{document}
